%!TEX root = ../template.tex
%%%%%%%%%%%%%%%%%%%%%%%%%%%%%%%%%%%%%%%%%%%%%%%%%%%%%%%%%%%%%%%%%%%%
%% appendix1.tex
%% NOVA thesis document file
%%
%% Appendix: step-by-step setup of the testbed, organised by machine
%%%%%%%%%%%%%%%%%%%%%%%%%%%%%%%%%%%%%%%%%%%%%%%%%%%%%%%%%%%%%%%%%%%%

\typeout{NT FILE appendix1.tex}%

%CHANGED BY CLAUDE: whole appendix restructured by machine (was "4G Network Setup").
% - The 4G introduction and the Hardware/Software Components sections were removed; the
%   hardware is described in Chapter 4, and the old text is in git history.
% - The 4G figures (4G architecture, 4G connection, IMG_9772 photo) were dropped in favour
%   of the testbed diagram and IP figure from Chapter 4.
% - Existing 4G steps are reused as-is unless marked "%CHANGED BY CLAUDE" below.
% - Sections marked TODO are placeholders for the 5G setup.

\chapter{Testbed Setup}
\label{app:testbed-setup}

This appendix is a step-by-step guide to rebuilding the testbed described in Chapter~\ref{cha:proposed-architecture}. It is organised by machine, following the layout in Figure~\ref{fig:testbed-architecture}: the steps in each section are run on the machine that section names, and a reader interested in only one generation can skip the section for the other.

\begin{itemize}
    \item \textbf{4G network:} Sections~\ref{sec:app-network}, \ref{sec:app-core}, \ref{sec:app-4g-ran} and~\ref{sec:app-running}.
    \item \textbf{5G network:} Sections~\ref{sec:app-network}, \ref{sec:app-core}, \ref{sec:app-5g-ran} and~\ref{sec:app-running}.
\end{itemize}

All machines run Ubuntu 22.04 LTS, and the software versions used in this dissertation are listed in Table~\ref{tab:software-versions}. The \gls{UE} additionally needs an eSIM profile, whose provisioning is described in Appendix~\ref{ch:provisioning_esim}.

%%--------------------------------------------------------------------------------------
%CHANGED BY CLAUDE (was "Connection between the machines"; now covers all machines through the switch)
\section{Network Connection (All Machines)}
\label{sec:app-network}

The machines are connected to each other through an Ethernet switch, each with a static IP address, as shown in Figure~\ref{fig:ip-addresses} and listed below. These connections carry the control and user plane traffic between the core and both \glspl{RAN}.

\begin{table}[H]
    \centering
    \begin{tabular}{|c|c|c|}
        \hline
        \textbf{Machine} & \textbf{Role} & \textbf{IP address} \\
        \hline
        1 & Core & 10.0.0.10/24 \\
        \hline
        2 & 5G \gls{RAN} & 10.0.0.20/24 \\
        \hline
        3 & 5G \gls{NSA} \gls{RAN} (unused) & 10.0.0.30/24 \\
        \hline
        4 & 4G \gls{RAN} & 10.0.0.40/24 \\
        \hline
    \end{tabular}
    \caption{Static IP addresses of the testbed machines.}
    \label{tab:app-ip-addresses}
\end{table}

The steps are the same on every machine, differing only in the IP address assigned. Connect the machine to the switch and check the available network interfaces using the following command:
\begin{code}
ip a
\end{code}
The output shows the available network interfaces, including the Ethernet interface (in this case enp0s31f6), as shown in Figure~\ref{fig:ip_a_output}.
\begin{figure}[H]
    \centering
    \includegraphics[width=\textwidth]{ip a.png} %CHANGED BY CLAUDE (placeholder replaced)
    \caption{Output of the \texttt{ip a} command on the core machine once set up, showing the Ethernet interface (enp0s31f6) at 10.0.0.10/24, and the \texttt{ogstun} interface created by Open5GS for the \gls{UE} address pool (10.45.0.1/16).}
    \label{fig:ip_a_output}
\end{figure}

Enable the Ethernet interface (if disabled) and assign a static IP address to it. On the core machine:

\begin{code}
sudo ip link set enp0s31f6 up
sudo ip addr add 10.0.0.10/24 dev enp0s31f6
\end{code}

To make the connection persistent across reboots, the network configuration needs to be added to the Netplan configuration file. This file may have a different name from the one in this example. Open the Netplan configuration file on \texttt{/etc/netplan/01-netcfg.yaml} using a text editor and add the following configuration to assign a static IP address to the Ethernet interface, replacing the address with the one for the machine being configured:
\begin{minted}[bgcolor=gray!20,frame=single,framerule=0.8pt,rulecolor=\color{black}]{yaml}
    # Let NetworkManager manage all devices on this system
    network:
      version: 2
      renderer: NetworkManager
      ethernets:
        enp0s31f6:
          dhcp4: false
          addresses:
            - 10.0.0.10/24
\end{minted}
After saving the changes, apply the new network configuration using the following command:
\begin{code}
sudo netplan apply
\end{code}

Finally, verify that the machines can reach each other by pinging each \gls{RAN} machine from the core machine:
\begin{code}
ping 10.0.0.20
ping 10.0.0.40
\end{code}
The output should show successful ping responses, as shown in Figure~\ref{fig:ping_output}.
\begin{figure}[H]
    \centering
    \includegraphics[width=0.85\textwidth]{pings.png} %CHANGED BY CLAUDE (placeholder replaced)
    \caption{Pings from the core machine to the 5G \gls{RAN} machine (10.0.0.20) and the 4G \gls{RAN} machine (10.0.0.40), with no packet loss and round-trip times of around 1\,ms.}
    \label{fig:ping_output}
\end{figure}

%%--------------------------------------------------------------------------------------
\section{Core Machine (Machine 1)}
\label{sec:app-core}

This section installs and configures Open5GS, which serves as the core for both generations. The installation process consists of the following steps:
\begin{enumerate}
    \item Installing MongoDB.
    \item Installing Open5GS.
    \item Installing the WebUI for Open5GS.
    \item Configuring the core network functions for the 4G and 5G networks.
    \item Routing the \gls{UE} traffic.
    \item Adding the subscribers.
\end{enumerate}
More information about the project and the installation steps can be found in the official Open5GS documentation~\cite{open5gs_docs}.

\subsection{Installing MongoDB}
\label{sec:Installing_MongoDB}
Open5GS depends on MongoDB as its database backend, where it stores subscriber information, network configurations, and operational data. Therefore, the first step in setting up Open5GS is to install MongoDB:
\begin{code}
sudo apt update
sudo apt install gnupg curl -y
curl -fsSL https://pgp.mongodb.com/server-8.0.asc | sudo gpg -o /usr/share/keyrings/mongodb-server-8.0.gpg --dearmor

echo "deb [ arch=amd64,arm64 signed-by=/usr/share/keyrings/mongodb-server-8.0.gpg ] https://repo.mongodb.com/apt/ubuntu jammy/mongodb-org/8.0 multiverse" | sudo tee /etc/apt/sources.list.d/mongodb-org-8.0.list

sudo apt update
sudo apt install -y mongodb-org
\end{code}
After executing the above commands, MongoDB will be installed on the machine. We can then start and verify the installation of the MongoDB daemon:
\begin{code}
sudo systemctl start mongod
sudo systemctl status mongod
\end{code}
In which the output, shown in Figure~\ref{fig:mongodb_status}, should indicate that the MongoDB service is active and running:
\begin{figure}[H]
    \centering
    \includegraphics[width=\textwidth]{4G_MongoDB_response.png}
    \caption{Output of the command to check MongoDB service status, showing that the service is active and running.}
    \label{fig:mongodb_status}
\end{figure}
To ensure MongoDB starts automatically on system boot, we can enable the service with the following command:
\begin{code}
sudo systemctl enable mongod
\end{code}
For more detailed instructions and troubleshooting guidance, refer to the official MongoDB installation documentation~\cite{mongodb_installation}.

\subsection{Installing Open5GS}
\label{sec:Installing_Open5GS}
After installing MongoDB, the next step is to add the Open5GS repository to the system and install Open5GS. This can be done by executing the following commands:
\begin{code}
sudo add-apt-repository ppa:open5gs/latest
sudo apt update
sudo apt install open5gs -y
\end{code}

This installs the latest Open5GS release available at the time; the testbed described in this dissertation ran version 2.8.0. After the installation is complete, we can verify that the Open5GS services are running correctly by checking their status:

\begin{code}
sudo systemctl status open5gs-*
\end{code}

\subsection{Installing the WebUI for Open5GS}
\label{sec:Installing_WebUI_for_Open5GS}
Open5GS provides a WebUI for easier management of subscriber information, network configurations, and monitoring. The interface depends on Node.js, which needs to be installed first (on most Ubuntu builds, Node.js comes already pre-installed).
To install Node.js, we can use the following commands:
\begin{code}
sudo apt update
sudo apt install -y ca-certificates curl gnupg
sudo mkdir -p /etc/apt/keyrings
curl -fsSL https://deb.nodesource.com/gpgkey/nodesource-repo.gpg.key |
sudo gpg --dearmor -o /etc/apt/keyrings/nodesource.gpg

NODE_MAJOR=20
echo "deb [signed-by=/etc/apt/keyrings/nodesource.gpg] https://deb.nodesource.com/node_$NODE_MAJOR.x nodistro main" |
sudo tee /etc/apt/sources.list.d/nodesource.list

sudo apt update
sudo apt install nodejs -y
\end{code}
After installing Node.js, we can proceed to install the Open5GS WebUI by running the following command:
\begin{code}
curl -fsSL https://open5gs.org/open5gs/assets/webui/install |
sudo -E bash -
\end{code}
The WebUI can be accessed by navigating to `http://localhost:9999` in a web browser. The default login credentials are:
\begin{itemize}
    \item Username: admin
    \item Password: 1423
\end{itemize}

\subsection{Configuring the Core Network Functions}
\label{sec:Configuring_Open5GS_Components}
In Open5GS, the configuration files for the network functions are located in the \texttt{/etc/open5gs/} directory, with each network function having its own file. The configuration files used in this dissertation are kept in the project's GitHub repository~\cite{thesis_repo}. The repository also holds a script, \texttt{deploy\_open5gs\_configs.sh}, which copies the files into that directory, keeping a backup of the existing ones, and restarts the Open5GS services:
\begin{code}
sudo ./deploy_open5gs_configs.sh
\end{code}
The following subsections describe the changes these files make to the Open5GS defaults, for each generation. Unless issued a PLMN ID by the national regulator, use the international test PLMN ID (001/01), which is the one used in this configuration.

\subsubsection{4G: MME and SGW-U}
Modify the MME configuration file, on \texttt{/etc/open5gs/mme.yaml}, to set the PLMN ID, TAC and the IP address the \gls{MME} listens on for the \gls{eNB}.

\begin{minted}{diff}
--- a/configs/open5gs/mme.yaml.in
+++ b/configs/open5gs/mme.yaml.in
@@ -10,7 +10,7 @@ mme:
     freeDiameter: @sysconfdir@/freeDiameter/mme.conf
     s1ap:
       server:
-        - address: 127.0.0.2
+        - address: 10.0.0.10
     gtpc:
       server:
         - address: 127.0.0.2
@@ -25,14 +25,14 @@ mme:
           port: 9090
     gummei:
       plmn_id:
-        mcc: 999
-        mnc: 70
+        mcc: 001
+        mnc: 01
       mme_gid: 2
       mme_code: 1
     tai:
       plmn_id:
-        mcc: 999
-        mnc: 70
+        mcc: 001
+        mnc: 01
-      tac: 1
+      tac: 7
     security:
         integrity_order : [ EIA2, EIA1, EIA0 ]
\end{minted}

Modify the SGW-U configuration file, on \texttt{/etc/open5gs/sgwu.yaml}, to set the GTP-U IP address.
\begin{minted}{diff}
--- a/configs/open5gs/sgwu.yaml.in
+++ b/configs/open5gs/sgwu.yaml.in
@@ -15,7 +15,7 @@ sgwu:
 #          - address: 127.0.0.3
     gtpu:
       server:
-        - address: 127.0.0.6
+        - address: 10.0.0.10
\end{minted}
After modifying the configuration files, restart the Open5GS daemons to apply all changes:
\begin{code}
sudo systemctl restart open5gs-mmed
sudo systemctl restart open5gs-sgwud
\end{code}

\subsubsection{5G: AMF, SMF, UPF and NRF}
Four 5G network functions need changes. The other network functions of the 5G core (AUSF, UDM, UDR, PCF, NSSF, BSF and SCP) keep their default configuration, since they only talk to each other over the loopback interface of the core machine.

Modify the AMF configuration file, on \texttt{/etc/open5gs/amf.yaml}, to set the IP address the \gls{AMF} listens on for the \gls{gNB}, and the PLMN ID and TAC, which must match the ones configured in the \gls{gNB}:
\begin{minted}{diff}
--- a/configs/open5gs/amf.yaml.in
+++ b/configs/open5gs/amf.yaml.in
@@ amf: @@
     ngap:
       server:
-        - address: 127.0.0.5
+        - address: 10.0.0.10
@@ amf: @@
     guami:
       - plmn_id:
-          mcc: 999
-          mnc: 70
+          mcc: 001
+          mnc: 01
         amf_id:
           region: 2
           set: 1
     tai:
       - plmn_id:
-          mcc: 999
-          mnc: 70
-        tac: 1
+          mcc: 001
+          mnc: 01
+        tac: 7
     plmn_support:
       - plmn_id:
-          mcc: 999
-          mnc: 70
+          mcc: 001
+          mnc: 01
         s_nssai:
           - sst: 1
\end{minted}

Modify the UPF configuration file, on \texttt{/etc/open5gs/upf.yaml}, to set the GTP-U address the \gls{gNB} sends user plane traffic to (the N3 interface):
\begin{minted}{diff}
--- a/configs/open5gs/upf.yaml.in
+++ b/configs/open5gs/upf.yaml.in
@@ upf: @@
     gtpu:
       server:
-        - address: 127.0.0.7
+        - address: 10.0.0.10
\end{minted}

Modify the SMF configuration file, on \texttt{/etc/open5gs/smf.yaml}, to name the data network of the \gls{UE} address pool and to set the DNS servers handed to the \glspl{UE}. Because the \gls{SMF} and \gls{UPF} also act as the PGW-C and PGW-U of the 4G network, these settings apply to both generations:
\begin{minted}{diff}
--- a/configs/open5gs/smf.yaml.in
+++ b/configs/open5gs/smf.yaml.in
@@ smf: @@
     session:
-      - subnet: 10.45.0.0/16
+      - dnn: internet
+        subnet: 10.45.0.0/16
         gateway: 10.45.0.1
       - subnet: 2001:db8:cafe::/48
         gateway: 2001:db8:cafe::1
     dns:
-      - 8.8.8.8
-      - 8.8.4.4
+      - 1.1.1.1
+      - 1.0.0.1
       - 2001:4860:4860::8888
       - 2001:4860:4860::8844
\end{minted}

Finally, modify the NRF configuration file, on \texttt{/etc/open5gs/nrf.yaml}, to set the PLMN it serves:
\begin{minted}{diff}
--- a/configs/open5gs/nrf.yaml.in
+++ b/configs/open5gs/nrf.yaml.in
@@ nrf: @@
     serving:
       - plmn_id:
-          mcc: 999
-          mnc: 70
+          mcc: 001
+          mnc: 01
\end{minted}
After modifying the configuration files, restart the Open5GS daemons to apply all changes:
\begin{code}
sudo systemctl restart open5gs-*
\end{code}

\subsection{Routing the UE Traffic}
\label{sec:app-routing}
For the \glspl{UE} to reach networks beyond the core machine, the core must forward their traffic and translate their addresses. Both generations assign \gls{UE} addresses from the same pool, so this is done once:
\begin{code}
### Enable IPv4/IPv6 Forwarding
$ sudo sysctl -w net.ipv4.ip_forward=1
$ sudo sysctl -w net.ipv6.conf.all.forwarding=1

### Add NAT Rule
$ sudo iptables -t nat -A POSTROUTING -s 10.45.0.0/16 ! -o ogstun -j MASQUERADE
$ sudo ip6tables -t nat -A POSTROUTING -s 2001:db8:cafe::/48 ! -o ogstun -j MASQUERADE
\end{code}

%CHANGED BY CLAUDE (placeholder)
\subsection{Adding the Subscribers}
\label{sec:app-subscribers}

Before a \gls{UE} can attach, its SIM must be registered in the core as a subscriber. Subscribers are added through the WebUI installed in Section~\ref{sec:Installing_WebUI_for_Open5GS}: open \texttt{http://localhost:9999} on the core machine, log in, and select Subscriber, which lists the subscribers already registered, as shown in Figure~\ref{fig:app-webui-list}.

\begin{figure}[H]
    \centering
    \includegraphics[width=\textwidth]{webUI.png}
    \caption{The Open5GS WebUI's subscriber list, each subscriber identified by its IMSI.}
    \label{fig:app-webui-list}
\end{figure}

The plus button at the bottom right opens the subscriber form, shown in Figure~\ref{fig:app-webui-add}. The fields must match the values programmed into the \gls{UE}'s SIM, which for the eSIM used in this project are given in Appendix~\ref{ch:provisioning_esim}:
\begin{itemize}
    \item \textbf{IMSI}: the subscriber's identity, whose first five digits (00101) are the PLMN of the testbed.
    \item \textbf{Subscriber Key (K)} and \textbf{Operator Key (OPc)}: the secret keys used to authenticate the SIM, with the USIM type set to OPc.
    \item \textbf{Authentication Management Field (AMF)}: left at its default, 8000.
\end{itemize}

The remaining fields, which set the slice and data network of the subscriber, are left at the WebUI's defaults (SST 1 and \texttt{internet}), which match the core configuration of Table~\ref{tab:core-config}. A single subscriber entry serves both generations: the same SIM authenticates against the 4G \gls{EPC} and the 5G core, since both read their subscriber data from the same database.

\begin{figure}[H]
    \centering
    \includegraphics[width=\textwidth]{webUI_add.png}
    \caption{The Open5GS WebUI's subscriber form, filled in for the testbed's \gls{UE}.}
    \label{fig:app-webui-add}
\end{figure}

\section{4G RAN Machine (Machine 4)}
\label{sec:app-4g-ran}

This section sets up the 4G \gls{RAN}: the drivers for the Nuand bladeRF xA4 and srsRAN 4G. The hardware itself, including the BT-200 and BT-100 amplifiers on the board's RX and TX ports, is described in Section~\ref{sec:hardware-components}.

\subsection{Setting up the Nuand drivers and firmware for the bladeRF xA4}
\label{sec:Nuand_BladeRF_xA4_FPGA_firmware}

Nuand provides official Ubuntu package repositories (PPAs) that simplify the installation of the bladeRF software in Ubuntu systems. Through these repositories, users can quickly install both the required host tools and libraries and the necessary firmware for the FPGA using \texttt{apt}, without needing to build everything manually from source. To enable the PPA:
\begin{code}
sudo add-apt-repository ppa:nuandllc/bladerf
sudo apt update
\end{code}
After enabling the PPA, we can install the bladeRF software and firmware using the following command:
\begin{code}
sudo apt install bladerf
sudo apt install bladerf-fpga-hostedx40
\end{code}

To bridge the connection between the bladeRF xA4 and the \gls{RAN}, SoapySDR will be used, which will be covered in the next section (\ref{sec:SoapySDR}). For this to work, the header files for the bladeRF drivers need to be installed:

\begin{code}
sudo apt install libbladerf-dev
\end{code}
With the bladeRF drivers installed, we can proceed into the installation of SoapySDR, which will allow us to use the bladeRF xA4 as the radio interface for the \gls{RAN}.

\subsection{Installing SoapySDR}
\label{sec:SoapySDR}

The SoapySDR component is comprised of two parts: the SoapySDR library, which is the main interface for all types of SDR devices, and the SoapyBladeRF plug-in, which provides specific support for the bladeRF xA4.

First, we need to install the necessary dependencies:
\begin{code}
sudo apt install -y cmake g++ libpython3-dev python3-numpy swig
\end{code}
Afterwards, we can clone the SoapySDR repository and build the library:
\begin{code}
git clone https://github.com/pothosware/SoapySDR.git
cd SoapySDR

mkdir build
cd build
cmake ..
make -j`nproc`
sudo make install -j`nproc`
sudo ldconfig
\end{code}
To verify the installation of the SoapySDR library, we can run the following command:
\begin{code}
SoapySDRUtil --info
\end{code}
After installing SoapySDR, we can proceed to install the SoapyBladeRF plug-in, through their repository:
\begin{code}
git clone https://github.com/pothosware/SoapyBladeRF.git
cd SoapyBladeRF

mkdir build
cd build
cmake ..
make
sudo make install
\end{code}
With both the SoapySDR library and the SoapyBladeRF plug-in installed, the bladeRF xA4 should now be recognized as an available SDR device, when running the following command:
\begin{code}
SoapySDRUtil --probe="driver=bladerf"
\end{code}
This command should output the details of the board, thus confirming that the bladeRF xA4 is properly configured and ready to be used as the radio interface for the 4G \gls{RAN}. With the SDR configuration complete, we can now proceed to install and configure srsRAN 4G.

\subsection{srsRAN 4G}
\label{sec:srsRAN_4G}
First, install the dependencies for srsRAN 4G:
\begin{code}
sudo apt-get install build-essential cmake libfftw3-dev libmbedtls-dev libboost-program-options-dev libconfig++-dev libsctp-dev
\end{code}
Then, clone the srsRAN repository and build the software:
\begin{code}
git clone https://github.com/srsRAN/srsRAN_4G.git
cd srsRAN_4G
mkdir build
cd build
cmake ../
make
make test
\end{code}
Proceed to install srsRAN 4G:
\begin{code}
sudo make install
srsran_install_configs.sh user
\end{code}

The second command installs the default configuration files in \texttt{\textasciitilde/.config/srsran/}. Only the \gls{eNB} files are used, since the core is provided by Open5GS rather than by srsRAN's own \gls{EPC}. The configuration used in this dissertation is kept in the project's repository~\cite{thesis_repo} (\texttt{enb.conf}, \texttt{rr.conf}, \texttt{sib.conf} and \texttt{rb.conf}). The lines that differ from the default files point the \gls{eNB} to the core, set its own address and set the radio parameters:
\begin{minted}[bgcolor=gray!20,frame=single,framerule=0.8pt,rulecolor=\color{black}]{ini}
# enb.conf
[enb]
enb_id = 0x19B
mcc = 001
mnc = 01
mme_addr = 10.0.0.10
gtp_bind_addr = 10.0.0.40
s1c_bind_addr = 10.0.0.40
n_prb = 50

[rf]
tx_gain = 66
rx_gain = 60
device_name = bladeRF

# rr.conf (cell list)
cell_id = 0x01;
tac = 0x0007;
pci = 1;
dl_earfcn = 3350;
\end{minted}

This completes the one-time setup of the 4G \gls{RAN} machine. The next steps to run the network are given in Section~\ref{sec:app-running}.

%%--------------------------------------------------------------------------------------

\section{5G RAN Machine (Machine 2)}
\label{sec:app-5g-ran}

This section sets up the 5G \gls{RAN}: the drivers for the NI B2901, the \gls{GPSDO} and the srsRAN Project \gls{gNB}. The hardware itself is described in Section~\ref{sec:hardware-components}.

\subsection{Installing UHD}
\label{sec:app-uhd}

The B2901 is driven through UHD, the USRP Hardware Driver. The testbed ran UHD version 4.10.0.0, built from source at that release rather than installed from the Ubuntu repositories, which may provide a different version. 
The OCUDU installation guide names UHD 4.6.0.0 as its requirement; the later 4.10.0.0 release worked with the \gls{gNB} used in this dissertation. First, install the build dependencies:
\begin{code}
sudo apt-get install autoconf automake build-essential ccache cmake cpufrequtils doxygen ethtool g++ git inetutils-tools libboost-all-dev libncurses-dev libusb-1.0-0 libusb-1.0-0-dev libusb-dev python3-dev python3-mako python3-numpy python3-requests python3-scipy python3-setuptools python3-ruamel.yaml libgrpc++-dev libprotobuf-dev protobuf-compiler-grpc
\end{code}
Then clone the UHD repository and check out the 4.10.0.0 release:
\begin{code}
git clone https://github.com/EttusResearch/uhd.git
cd uhd
git checkout v4.10.0.0
\end{code}
The UHD host software is built from the \texttt{host} directory of the repository:
\begin{code}
cd host
mkdir build
cd build
cmake ../
make
make test
sudo make install
sudo ldconfig
\end{code}
The B2901 loads its FPGA image from the host each time it is powered on, so the images matching the installed UHD version must be downloaded:
\begin{code}
sudo uhd_images_downloader
\end{code}
Finally, with the B2901 connected over USB 3.0, check that UHD detects it, and that it reports the board's details without errors:
\begin{code}
uhd_find_devices
uhd_usrp_probe
\end{code}
The second command loads the FPGA image onto the board, so it takes a few seconds to complete the first time it is run after the board is powered on. 

\subsection{Setting up the GPSDO}
\label{sec:app-gpsdo}

The B2901 synchronizes its frequency and time references to the Leo Bodnar \gls{GPSDO}, which in turn disciplines its oscillator to the GPS satellites. The \gls{GPSDO} provides two signals, each connected to the matching input on the rear of the B2901, as shown in Figure~\ref{fig:b2901-gpsdo}:
\begin{itemize}
  \item a 1PPS (one pulse per second) signal, on Output 1, connected to \texttt{PPS IN} port, which aligns the board's time to the start of each second.
    \item a 10\,MHz reference clock, on Output 2, connected to \texttt{REF IN} port, from which the board derives all of its sampling and carrier frequencies;
\end{itemize}

The \gls{gNB} YAML file is configured to use both through the \texttt{clock: external} and \texttt{sync: external} settings (Section~\ref{sec:app-gnb-config}). The \gls{GPSDO} is powered over USB from Machine 2, and the same USB connection is used to configure it and to check its lock.

\subsubsection{Placing the GPS antenna}
The \gls{GPSDO} can only discipline its oscillator while it receives signal from the GPS satellites, so its antenna needs a clear view of the sky. In this testbed, the antenna was placed on the windowsill of the lab, as shown in Figure~\ref{fig:app-gps-antenna}. After power-on, the receiver locks the signal to the satellites after a few minutes.

\begin{figure}[H]
    \centering
    \includegraphics[width=\textwidth]{GPSDO_Windowsill.jpeg} %CHANGED BY CLAUDE (placeholder replaced with the photo)
    \caption{The \gls{GPSDO}'s GPS antenna, placed on the windowsill outside the lab window for a clear view of the sky.}
    \label{fig:app-gps-antenna}
\end{figure}

\subsubsection{Installing lbgpsdo}
The \gls{GPSDO} is configured and monitored with \texttt{lbgpsdo}~\cite{lbgpsdo}, an open-source command-line utility for Leo Bodnar devices. It is written in Python and depends on the \texttt{hid} package, so it is installed in its own virtual environment:

\begin{code}
git clone https://github.com/hamarituc/lbgpsdo.git
cd lbgpsdo
python3 -m venv venv
source venv/bin/activate
pip install -r requirements.txt
\end{code}

The utility provides two scripts: \texttt{lbe142x.py}, for the LBE-1421 used in this testbed, and \texttt{lbgpsdo.py}, for the older Leo Bodnar GPSDO models. With the \gls{GPSDO} connected, confirm that it is detected:

\begin{code}
./lbe142x.py list
\end{code}

\subsubsection{Configuring the outputs}
%CHANGED BY CLAUDE (command matched to the output assignment above: Output 1 = 1PPS, Output 2 = 10 MHz; TODO removed)
Set Output 1 to the 1PPS signal and Output 2 to 10\,MHz, and save the settings to the device, so that they persist when it is powered off:
\begin{code}
./lbe142x.py modify --pps-enable --f2 10000000 --save
\end{code}

\subsubsection{Checking the lock}
Before starting the \gls{gNB}, check that the \gls{GPSDO} is locked:
\begin{code}
sudo ./lbe142x.py detail
\end{code}
%CHANGED BY CLAUDE (sudo added to match the screenshot; TODO replaced with the screenshot)
The output shows the device's status and the settings of each output, as shown in Figure~\ref{fig:app-lbe-detail}. The \gls{GPSDO} is ready once it reports both satellite lock (\texttt{SAT lock}), meaning it is receiving enough satellites, and PLL lock, meaning its output is disciplined to them. The output settings confirm the configuration from the previous step: 1PPS on Output 1 and 10\,MHz on Output 2.

\begin{figure}[H]
    \centering
    \includegraphics[width=0.6\textwidth]{lbe142x.png}
    \caption{Output of \texttt{lbe142x.py detail} for the LBE-1421, reporting both satellite and PLL lock, and the two outputs set to 1PPS and 10\,MHz.}
    \label{fig:app-lbe-detail}
\end{figure}

\subsection{srsRAN Project (OCUDU)}
\label{sec:app-srsran-project}

The srsRAN Project \gls{gNB} is now developed under the name OCUDU~\cite{ocudu_project}.
UHD must be installed, as described in Section~\ref{sec:app-uhd}, before the \gls{gNB} is built. The build checks for UHD only when it is configured, and includes support for the B2901 only if UHD is found at that point. 

A \gls{gNB} built without UHD pre-installed compiles and installs without errors, but will not  detect the radio when started, and installing UHD afterwards does not fix it. The \gls{gNB} has to be rebuilt from the \texttt{cmake} step onwards. With UHD in place, install the build dependencies:

\begin{code}
sudo apt-get install cmake make gcc g++ pkg-config libmbedtls-dev libsctp-dev libyaml-cpp-dev libgtest-dev
\end{code}

Then clone the repository and check out the commit used in this dissertation, on the project's \texttt{dev} branch: %CHANGED BY CLAUDE (checkout restored; commit confirmed by git log on Machine 2)

\begin{code}
git clone https://gitlab.com/ocudu/ocudu.git
cd ocudu
\end{code}

Build the software, run its tests, and install it:
\begin{code}
mkdir build
cd build
cmake ../
make -j $(nproc)
make test -j $(nproc)
sudo make install
\end{code}

The build takes several minutes. Once installed, the \gls{gNB} is available system-wide as the \texttt{gnb} command, which is used to start the 5G \gls{RAN} in Section~\ref{sec:app-start-5g}.

\subsection{Configuring the gNB}
\label{sec:app-gnb-config}

Unlike srsRAN 4G, the srsRAN Project \gls{gNB} reads a single YAML configuration file, passed to it through the startup command. The file used in this dissertation is kept in the project's repository~\cite{thesis_repo} as \texttt{gnb.yaml}, and is shown below, as used in Round 3. It has five parts:
\begin{itemize}
    \item \texttt{cu\_cp}: the connection to the \gls{AMF} on the core machine, and the tracking area, PLMN and slice the \gls{gNB} supports, which must match the core configuration.
    \item \texttt{ru\_sdr}: the radio. The B2901 is driven through UHD, takes both its clock and its time reference from the \gls{GPSDO} (\texttt{external}), and runs at the transmit and receive gains discussed in Chapter~\ref{cha:methodology}.
    \item \texttt{cell\_cfg}: the cell itself, in band n78 with a 20\,MHz channel and 30\,kHz subcarrier spacing, using the 2:2 TDD pattern.
    \item \texttt{log} and \texttt{metrics}: logging of all layers at \texttt{info} level, with the per-second metrics written into the log, as described in Section~\ref{sec:data-collection}.
    \item \texttt{pcap}: packet captures of the MAC layer and of the \gls{NGAP} signalling with the core.
\end{itemize}
\begin{minted}[bgcolor=gray!20,frame=single,framerule=0.8pt,rulecolor=\color{black}]{yaml}
cu_cp:
  inactivity_timer: 7200
  amf:
    addrs: 10.0.0.10
    port: 38412
    bind_addrs: 10.0.0.20
    supported_tracking_areas:
      - tac: 7
        plmn_list:
          - plmn: "00101"
            tai_slice_support_list:
              - sst: 1

ru_sdr:
  device_driver: uhd
  device_args: type=b200
  srate: 23.04
  clock: external
  sync: external
  otw_format: sc12
  tx_gain: 89.75
  rx_gain: 60

cell_cfg:
  dl_arfcn: 627340
  band: 78
  channel_bandwidth_MHz: 20
  common_scs: 30
  plmn: "00101"
  tac: 7
  pci: 1
  tdd_ul_dl_cfg:
    dl_ul_tx_period: 5
    nof_dl_slots: 2
    nof_dl_symbols: 5
    nof_ul_slots: 2
    nof_ul_symbols: 5
  q_rx_lev_min: -70

log:
  filename: /tmp/gnb.log
  all_level: info
  hex_max_size: 0

metrics:
  autostart_stdout_metrics: true
  enable_log: true
  enable_verbose: true

pcap:
  mac_enable: true
  mac_filename: /tmp/gnb_mac.pcap
  ngap_enable: true
  ngap_filename: /tmp/gnb_ngap.pcap
\end{minted}
Note that the log and packet captures are written to \texttt{/tmp}, which is cleared when the machine restarts, so they must be copied elsewhere after each test session.

This completes the one-time setup of the 5G \gls{RAN} machine. The steps repeated each time the network is started are given in Section~\ref{sec:app-running}.

%%--------------------------------------------------------------------------------------
\section{Running the Network}
\label{sec:app-running}

The previous sections cover setup, which is done once. This section covers the steps repeated every time the network is started. They assume the core has been configured and the subscriber added (Section~\ref{sec:app-core}); since the Open5GS network functions run as system services, the core is already running whenever the core machine is on.

\subsection{Preparing the RAN Machines}
To ensure the best performance on the \gls{RAN} side, the srsRAN team provides a performance script to optimise the system for real-time processing. This script should be executed on each \gls{RAN} machine before starting its \gls{RAN}:
\begin{code}
# Download and execute the srsRAN performance tuning script
BASE=https://raw.githubusercontent.com/srsran/srsRAN_Project/main/scripts
curl -fsSL $BASE/srsran_performance | bash
\end{code}


\subsection{Starting the 4G Network}
\label{sec:app-start-4g}
To start the 4G network, start the \gls{eNB} on Machine 4: 
\begin{code}
sudo srsenb
\end{code}
The \gls{eNB} connects to the \gls{MME} at 10.0.0.10. Once it reports the S1 connection as established, the 4G \gls{RAN} is up and running, and the \gls{UE} can connect to the network.

\subsection{Starting the 5G Network}
\label{sec:app-start-5g}
%CHANGED BY CLAUDE (drafted; the expected output is taken from the Round 3 gnb1.log and amf.log)
Before starting the 5G network, confirm that the \gls{GPSDO} has GPS lock (Section~\ref{sec:app-gpsdo}), as the B2901 takes both its clock and its time reference from it, and the \gls{gNB} will not transmit correctly without them. Then start the \gls{gNB} on Machine 2, passing it the configuration file from Section~\ref{sec:app-gnb-config}:
\begin{code}
sudo gnb -c gnb.yaml
\end{code}
On start-up, the \gls{gNB} opens an \gls{SCTP} connection to the \gls{AMF} at 10.0.0.10 and exchanges an NG Setup Request and Response with it. The following lines in \texttt{/tmp/gnb.log} show a successful start:
\begin{code}
[CU-CP] N2: Connection to AMF on 10.0.0.10:38412 was established
[NGAP ] Tx PDU: NGSetupRequest
[NGAP ] Rx PDU: NGSetupResponse
[CU-CP] Connected to AMF. Supported PLMNs: 00101
[GNB  ] ==== gNB started ===
\end{code}
The same exchange can be confirmed on the core machine, where the \gls{AMF} reports the \gls{gNB} as connected in its log file, \texttt{/var/log/open5gs/amf.log}:
\begin{code}
[amf] INFO: gNB-N2 accepted[10.0.0.20] in master_sm module
[amf] INFO: [Added] Number of gNBs is now 1
\end{code}
Once the \gls{gNB} has started, the 5G \gls{RAN} is up and running, and the \gls{UE} can connect to the network. With the configuration used in this dissertation, the \gls{gNB} also prints its per-second metrics to the console as soon as a \gls{UE} is attached. To stop the \gls{gNB}, press \texttt{Ctrl+C}. 

The log and packet captures in \texttt{/tmp} should then be copied elsewhere, as noted in Section~\ref{sec:app-gnb-config}.

\subsection{Verifying the Connection}
\textcolor{red}{TODO:
IMSI attach from core, network shows up on phone.
Start iperf3 server on core.
Cellular Lab running the iperf3 test.}

\section{Troubleshooting}
\label{sec:Troubleshooting}
\textcolor{red}{TODO.}
% Candidates from the test campaigns: receiver overload at short range (rx_gain), the TDD
% pattern capping uplink airtime, GPSDO not locked, RLC queue bufferbloat.
% OCUDU rebuild after UHD install, etc.